\chapter{Influenza Testing}

\section*{What Is This Test?}
Influenza testing detects influenza A and B viruses, the causative agents of seasonal epidemic influenza. Influenza viruses are enveloped, single-stranded RNA orthomyxoviruses. Influenza A is subtyped by hemagglutinin (H1-H18) and neuraminidase (N1-N11) surface glycoproteins; currently circulating subtypes include H1N1 and H3N2. Influenza B has two lineages: Victoria and Yamagata. These viruses cause acute febrile respiratory illness with abrupt onset of fever, myalgia, headache, sore throat, and nonproductive cough, typically occurring in winter epidemics. Complications include primary viral pneumonia, secondary bacterial pneumonia (most commonly \textit{Staphylococcus aureus} and \textit{Streptococcus pneumoniae}), exacerbation of underlying chronic disease, myositis, myocarditis, and neurologic complications including encephalitis and Guillain-Barré syndrome. High-risk groups for severe disease include children under 5 years, adults over 65, pregnant women, and individuals with chronic medical conditions.

\section*{Alternative Names}
Flu Test; Influenza A/B PCR; Rapid Influenza Diagnostic Test (RIDT); Influenza Antigen Test; Influenza Culture; Influenza DFA

\section*{Principle}
Multiple detection methods are available. Rapid influenza diagnostic tests (RIDTs) detect influenza A and B nucleoprotein antigens by lateral-flow immunochromatographic assay with results in 10--15 minutes \cite{merckx2017diagnostic}. Direct fluorescent antibody (DFA) staining uses fluorescein-labeled monoclonal antibodies to detect viral antigens in respiratory epithelial cells. Reverse-transcription polymerase chain reaction (RT-PCR) on nasopharyngeal specimens detects influenza virus RNA with the highest sensitivity (90--95\%) and specificity (95--98\%) and is the gold standard. RT-PCR can subtype influenza A viruses and distinguish influenza B lineages. Rapid molecular assays (e.g., Roche Cobas Liat, Cepheid GeneXpert) use isothermal nucleic acid amplification or cartridge-based RT-PCR \cite{vos2019rapid}. Viral culture in primary rhesus monkey kidney or MDCK cells is the reference method for virus isolation and subsequent antigenic characterization but requires 3--7 days. Serologic testing by hemagglutination inhibition (HAI) or microneutralization assays measures anti-hemagglutinin antibodies for epidemiologic and vaccine studies.

\section*{History}
Influenza viruses were first isolated in 1933 by Wilson Smith, Christopher Andrewes, and Patrick Laidlaw, who infected ferrets with filtered nasal washings from patients with influenza. The first influenza B virus was isolated by Thomas Francis Jr. in 1940. The phenomenon of hemagglutination was discovered by George Hirst in 1941, providing the basis for the HAI serologic assay. Viral culture techniques in embryonated eggs and then in tissue culture allowed large-scale vaccine production. Direct and indirect fluorescent antibody tests for rapid diagnosis were developed in the 1960s--1970s. The first rapid antigen tests (EIA-based) became available in the 1990s. RT-PCR for influenza was introduced in the late 1990s and became widely available after the 2009 H1N1 pandemic. Multiplex molecular panels capable of detecting influenza alongside other respiratory pathogens (e.g., BioFire FilmArray) were FDA-cleared starting in 2011. Point-of-care molecular tests were introduced in the 2010s for rapid bedside diagnosis.

\section*{Why This Test Is Ordered}
\begin{itemize}
  \item Rapid diagnosis of influenza A or B infection in patients presenting with acute febrile respiratory illness during influenza season (typically November through April in temperate regions of the Northern Hemisphere)
  \item Initiation of antiviral therapy (oseltamivir, zanamivir, baloxavir marboxil) within 48 hours of symptom onset to reduce illness duration and complications \cite{uyeki2019clinical}
  \item Differentiation of influenza from other respiratory viral infections (RSV, parainfluenza, adenovirus, human metapneumovirus) with overlapping clinical presentations
  \item Infection control decisions for hospitalized patients, including initiation of droplet precautions and cohorting
  \item Outbreak investigation in closed settings (nursing homes, prisons, dormitories, cruise ships)
  \item Surveillance monitoring for circulating influenza strains and assessment of vaccine match, coordinated by the CDC and WHO Global Influenza Surveillance and Response System (GISRS) \cite{cdc2022influenza}
  \item Evaluation of patients with influenza-like illness (ILI), defined as fever 100\textdegree{}F or greater plus cough and/or sore throat, in the absence of another known cause
\end{itemize}

\section*{Is This a Primary or Secondary Test}
This is a primary diagnostic test for suspected influenza infection. During periods of high influenza activity, clinical diagnosis alone may suffice. RT-PCR is preferred for hospitalized patients and those at high risk for complications. Rapid antigen testing may be used for initial screening in outpatient settings, but negative results should be confirmed with RT-PCR when clinical suspicion is high given the limited sensitivity (50--70\%) of RIDTs.

\section*{How This Test Is Performed}
Specimen collection: Nasopharyngeal (NP) swab is the preferred specimen. A flexible swab with a flocked nylon tip is inserted into one nostril to the nasopharynx (distance equal to that from the ear to the nostril), rotated against the mucosa for 5--10 seconds, then removed and placed in viral transport medium (VTM). Alternative specimens include nasal aspirate, nasal wash, throat swab (lower sensitivity), combined NP/throat swab, endotracheal aspirate (in intubated patients), or bronchoalveolar lavage (BAL). For intubated patients, lower respiratory tract specimens may have higher yield. Specimens should be collected within 72 hours of symptom onset for optimal sensitivity. Specimens should be refrigerated at 2--8\textdegree{}C and transported to the laboratory within 72 hours, or frozen at -70\textdegree{}C for longer storage.

\section*{What Preparation Is Needed}
No special patient preparation is required. Patients should avoid eating or drinking 30 minutes before specimen collection if an oropharyngeal swab is used. Prior use of intranasal medications or saline sprays should be noted as they may reduce viral recovery. Antiviral therapy started before specimen collection may reduce viral load and decrease test sensitivity.

\section*{Contraindications and Precautions}
Nasopharyngeal swab collection may cause transient discomfort, gagging, or sneezing. Relative contraindications include severe coagulopathy, recent nasal surgery, or severe nasal obstruction. False-negative rapid antigen test results are common (sensitivity 50--70\%), particularly when viral loads are low, in early or late infection, or with certain strains. Recent administration of live attenuated influenza vaccine (LAIV) within 7 days can cause false-positive results in antigen- and PCR-based tests. Hemolyzed or lipemic specimens may interfere with serologic assays but not with molecular tests.

\section*{Risks}
Nasopharyngeal swabbing may cause minor discomfort, transient gagging, sneezing, or mild epistaxis. There is essentially no risk from the collection procedure beyond these minor effects. All specimens should be collected using standard precautions and appropriate personal protective equipment given the potential for aerosol generation.

\section*{What Happens After the Test}
Rapid antigen test results are available within 10--15 minutes. Rapid molecular test results are available within 15--30 minutes. Standard RT-PCR results are typically available within 2--4 hours. Viral culture results may take 3--7 days. Positive results should prompt initiation of antiviral therapy and appropriate infection control measures. For hospitalized patients, negative RIDT results should be confirmed by molecular testing. Oseltamivir treatment should not be delayed while awaiting test results in high-risk patients during influenza season \cite{ison2019therapeutic}.

\section*{Understanding Results}

\subsection*{Below Normal}
\begin{itemize}
  \item A negative influenza test (antigen negative, PCR not detected) means influenza virus RNA or antigen was not detected in the specimen. This may represent: True absence of influenza infection in an uninfected individual; a viral infection caused by another respiratory pathogen (RSV, rhinovirus, adenovirus); or a false-negative result due to inadequate specimen collection, testing outside the window of peak viral shedding (before 24 hours or after 5--7 days from symptom onset), or use of a low-sensitivity rapid antigen test
  \item A negative rapid antigen test in a patient with high clinical suspicion for influenza (during seasonal epidemic, typical symptoms) should be confirmed by RT-PCR, as the sensitivity of RIDTs ranges from 50--70\% (high rate of false negatives)
\end{itemize}

\subsection*{Above Normal}
\begin{itemize}
  \item Influenza A detected: Positive RT-PCR or antigen test for influenza A indicates active influenza A virus infection. Clinical manifestations range from mild upper respiratory illness to severe viral pneumonia with acute respiratory distress syndrome (ARDS). Results may include subtype information (e.g., H1N1, H3N2). Certain subtypes (e.g., avian influenza A H5N1, H7N9) are associated with higher mortality and require public health notification
  \item Influenza B detected: Positive RT-PCR or antigen test for influenza B indicates active influenza B virus infection. Generally associated with milder illness than influenza A in adults but can cause severe disease in children. Influenza B is not subtyped by hemagglutinin and neuraminidase but may be reported as Victoria or Yamagata lineage
  \item Co-detection of influenza with bacterial pathogens (particularly \textit{S. aureus, S. pneumoniae, S. pyogenes}) on respiratory culture suggests secondary bacterial pneumonia, a leading cause of influenza-associated mortality
\end{itemize}

\section*{Normal Results}
\textbf{Negative} or \textbf{Not detected} by RT-PCR, DFA, or rapid antigen testing. For serology: A fourfold or greater rise in HAI antibody titer between acute (collected within 7 days of illness onset) and convalescent (collected 2--4 weeks later) serum samples is diagnostic of recent infection; a single convalescent titer of 1:32 or greater is suggestive but not definitive.

\section*{Follow-up Tests}
\begin{itemize}
  \item \textbf{Influenza A/B RT-PCR with subtyping:} Confirmatory testing when rapid antigen test is negative with high clinical suspicion; provides subtype information (H1, H3, H1N1) relevant to antiviral resistance patterns and public health surveillance
  \item \textbf{Respiratory pathogens panel (multiplex PCR):} Simultaneous detection of 15--20+ respiratory pathogens (viral and bacterial) when alternative or co-infecting pathogens are suspected
  \item \textbf{Chest radiograph:} To evaluate for viral or secondary bacterial pneumonia when respiratory symptoms are severe or hypoxemia is present
  \item \textbf{Sputum culture and blood cultures:} When secondary bacterial pneumonia or bacteremia is suspected, especially in patients with persistent or worsening symptoms after initial improvement
  \item \textbf{Viral culture:} For isolation and characterization of circulating influenza strains for vaccine strain selection and antiviral resistance surveillance (primarily performed at reference laboratories)
\end{itemize}

\section*{References}
\bibliographystyle{unsrtnat}
\bibliography{references}

\clearpage
